                       IMC2007, Blagoevgrad, Bulgaria
                                   Day 2, August 6, 2007

Problem 1. Let f : R → R be a continuous function. Suppose that for any c > 0, the graph
of f can be moved to the graph of cf using only a translation or a rotation. Does this imply that
f (x) = ax + b for some real numbers a and b ?
Solution. No. The function f (x) = ex also has this property since cex = ex+log c .
Problem 2. Let x, y, and z be integers such that S = x4 + y 4 + z 4 is divisible by 29. Show that S
is divisible by 294 .
Solution. We claim that 29 | x, y, z. Then, x4 + y 4 + z 4 is clearly divisible by 294 .
    Assume, to the contrary, that 29 does not divide all of the numbers x, y, z. Without loss of
generality, we can suppose that 29 ∤ x. Since the residue classes modulo 29 form a field, there is some
w ∈ Z such that xw ≡ 1 (mod 29). Then, (xw)4 + (yw)4 + (zw)4 is also divisible by 29. So we can
assume that x ≡ 1 (mod 29).
    Thus, we need to show that y 4 + z 4 ≡ −1 (mod 29), i.e. y 4 ≡ −1 − z 4 (mod 29), is impossible.
There are only eight fourth powers modulo 29,

                              0   ≡   04 ,
                              1   ≡   14 ≡ 124 ≡ 174 ≡ 284 (mod 29),
                              7   ≡   84 ≡ 94 ≡ 204 ≡ 214 (mod 29),
                             16   ≡   24 ≡ 54 ≡ 244 ≡ 274 (mod 29),
                             20   ≡   64 ≡ 144 ≡ 154 ≡ 234 (mod 29),
                             23   ≡   34 ≡ 74 ≡ 224 ≡ 264 (mod 29),
                             24   ≡   44 ≡ 104 ≡ 194 ≡ 254 (mod 29),
                             25   ≡   114 ≡ 134 ≡ 164 ≡ 184 (mod 29).

The differences −1 − z 4 are congruent to 28, 27, 21, 12, 8, 5, 4, and 3. None of these residue classes
is listed among the fourth powers.
Problem 3. Let C be a nonempty closed bounded subset of the real line and f : C → C be a
nondecreasing continuous function. Show that there exists a point p ∈ C such that f (p) = p.
   (A set is closed if its complement is a union of open intervals. A function g is nondecreasing if
g(x) ≤ g(y) for all x ≤ y.)
Solution. Suppose f (x) 6= x for all x ∈ C. Let [a, b] be the smallest closed interval that contains C.
Since C is closed, a, b ∈ C. By our hypothesis f (a) > a and f (b) < b. Let p = sup{x ∈ C : f (x) > x}.
Since C is closed and f is continuous, f (p) ≥ p, so f (p) > p. For all x > p, x ∈ C we have f (x) < x.
Therefore f f (p) < f (p) contrary to the fact that f is non-decreasing.
Problem 4. Let n > 1 be an odd positive integer and A = (aij )i,j=1...n be the n × n matrix with
                                 
                                 2 if i = j
                                 
                            aij = 1 if i − j ≡ ±2 (mod n)
                                 
                                 
                                   0 otherwise.

Find det A.



                                                  1
                                                  
                                  2                   1 if i − j ≡ ±1 (mod n)
Solution. Notice that A = B , with bij =                                      . So it is sufficient to find
                                                      0         otherwise
det B.
    To find det B, expand the determinant with respect to the first row, and then expad both terms
with respect to the first column.

              0 1                     1
                                                1 1              1                0 1
              1 0 1
                                                  01                              1 0  1
                1 0   1                           .    .                              .    .
                    .. ..                     1 .. ..                               1 .. ..
 det B =       1        .     .       =−                       +
                                                  ..                                  ..
                    ..                               . 0 1                               . 0 1
                        . 0 1
                                                        1 0 1                               1 0
                             1 0 1
                                          1                1 0   1                            1
         1                       1 0
                                                                                                        
             0 1                      1                        1 0                 1          0 1
                .. ..                0   1                                     .. ..       1 0 1        
           1       .      .                                   1                   .  .                  
                ..                      .. ..                                  ..              .. ..    
       = −         . 0 1           − 1     .   .          +
                                                                                   . 0 1  −   1    .  .  
                                        ..                                                     ..       
                         1 0 1             . 0 1                                    1 0           . 0 1
                                1 0            1 0 1                                      1            1 0
       = −(0 − 1) + (1 − 0) = 2,

since the second and the third matrices are lower/upper triangular, while in the first and the fourth
matrices we have row1 − row3 + row5 − · · · ± rown−2 = 0̄.
    So det B = 2 and thus det A = 4.
Problem 5. For each positive integer k, find the smallest number nk for which there exist real
nk × nk matrices A1 , A2 , . . . , Ak such that all of the following conditions hold:

 (1) A21 = A22 = . . . = A2k = 0,

 (2) Ai Aj = Aj Ai for all 1 ≤ i, j ≤ k, and

 (3) A1 A2 . . . Ak 6= 0.

Solution. The anwser is nk = 2k . In that case, the matrices can be constructed as follows: Let V be
the n-dimensional real vector space with basis elements [S], where S runs through all n = 2k subsets
of {1, 2, . . . , k}. Define Ai as an endomorphism of V by
                                                  (
                                                   0         if i ∈ S
                                         Ai [S] =
                                                   [S ∪ {i}] if i 6∈ S

for all i = 1, 2, . . . , k and S ⊂ {1, 2, . . . , k}. Then A2i = 0 and Ai Aj = Aj Ai . Furthermore,

                                      A1 A2 . . . Ak [∅] = [{1, 2, . . . , k}],

and hence A1 A2 . . . Ak 6= 0.
    Now let A1 , A2 , . . . , Ak be n × n matrices satisfying the conditions of the problem; we prove that
n ≥ 2k . Let v be a real vector satisfying A1 A2 . . . Ak v 6= 0. Denote by P the set of all subsets of
{1, 2, . . . , k}. Choose a complete ordering ≺ on P with the property

                               X≺Y        ⇒      |X| ≤ |Y | for all X, Y ∈ P.

                                                         2
For every element X = {x1 , x2 , . . . , xr } ∈ P, define AX = Ax1 Ax2 . . . Axr and vX = AX v. Finally,
write X̄ = {1, 2, . . . , k} \ X for the complement of X.
    Now take X, Y ∈ P with X  Y . Then AX̄ annihilates vY , because X  Y implies the existence
of some y ∈ Y \ X = Y ∩ X̄, and

                                     AX̄ vY = AX̄\{y} Ay Ay vY \{y} = 0,

since A2y = 0. So, AX̄ annihilates the span of all the vY with X  Y . This implies that vX does not
lie in this span, because AX̄ vX = v{1,2,...,k} 6= 0. Therefore, the vectors vX (with X ∈ P) are linearly
independent; hence n ≥ |P| = 2k .
Problem 6. Let f 6= 0 be a polynomial with real coefficients. Define the sequence f0 , f1 , f2 , . . . of
polynomials by f0 = f and fn+1 = fn + fn′ for every n ≥ 0. Prove that there exists a number N such
that for every n ≥ N, all roots of fn are real.
Solution. For the proof, we need the following
    Lemma 1. For any polynomial g, denote by d(g) the minimum distance of any two of its real
zeros (d(g) = ∞ if g has at most one real zero). Assume that g and g + g ′ both are of degree k ≥ 2
and have k distinct real zeros. Then d(g + g ′) ≥ d(g).
    Proof of Lemma 1: Let x1 < x2 < · · · < xk be the roots of g. Suppose a, b are roots of g + g ′
satisfying 0 < b − a < d(g). Then, a, b cannot be roots of g, and

                                              g ′ (a)   g ′ (b)
                                                      =         = −1.                                    (1)
                                              g(a)      g(b)
       ′
Since gg is strictly decreasing between consecutive zeros of g, we must have a < xj < b for some j.
    For all i = 1, 2, . . . , k − 1 we have xi+1 − xi > b − a, hence a − xi > b − xi+1 . If i < j, both sides
                                                                                    1
of this inequality are negative; if i ≥ j, both sides are positive. In any case, a−x  i
                                                                                        < b−x1i+1 , and hence

                              k−1                 k−1
                     g ′ (a) X 1            1     X      1        1      g ′ (b)
                            =           +       <             +        =
                     g(a)     i=1
                                  a − xi a − xk       b − xi+1 b − x1    g(b)
                                          | {z } i=1            | {z }
                                                 <0                         >0

This contradicts (1).
     Now we turn to the proof of the stated problem. Denote by m the degree of f . We will prove
by induction on m that fn has m distinct real zeros for sufficiently large n. The cases m = 0, 1 are
trivial; so we assume m ≥ 2. Without loss of generality we can assume that f is monic. By induction,
the result holds for f ′ , and by ignoring the first few terms we can assume that fn′ has m − 1 distinct
                                                           (n)   (n)           (n)
real zeros for all n. Let us denote these zeros by x1 > x2 > · · · > xm−1 . Then fn has minima
       (n)  (n)  (n)                         (n)  (n)  (n)                                  (n)   (n)
in x1 , x3 , x5 , . . . , and maxima in x2 , x4 , x6 , . . . . Note that in the interval (xi+1 , xi ), the
              ′
function fn+1     = fn′ + fn′′ must have a zero (this follows by applying Rolle’s theorem to the function
                                                       (n)
ex fn′ (x)); the same is true for the interval (−∞, xm−1 ). Hence, in each of these m − 1 intervals, fn+1
                                                                                                       ′

has exactly one zero. This shows that
                            (n)     (n+1)       (n)     (n+1)     (n)   (n+1)
                           x1 > x1          > x2 > x2           > x3 > x3       > ...                    (2)

                                        (n)                                        (n)
   Lemma 2. We have limn→∞ fn (xj ) = −∞ if j is odd, and lim fn (xj ) = +∞ if j is even.
                                                                        n→∞
    Lemma 2 immediately implies the result: For sufficiently large n, the values of all maxima of fn
are positive, and the values of all minima of fn are negative; this implies that fn has m distinct zeros.


                                                        3
  Proof of Lemma 2: Let d = min{d(f ′), 1}; then by Lemma 1, d(fn′ ) ≥ d for all n. Define
   (m − 1)dm−1
ε=             ; we will show that
      mm−1
                                          (n+1)                (n)
                                  fn+1 (xj        ) ≥ fn (xj ) + ε for j even.                                      (3)
                                                                                                              (n)
(The corresponding result for odd j can be shown similarly.) Do to so, write f = fn , b = xj , and
choose a satisfying d ≤ b − a ≤ 1 such that f ′ has no zero inside (a, b). Define ξ by the relation
        1
b − ξ = (b − a); then ξ ∈ (a, b). We show that f (ξ) + f ′ (ξ) ≥ f (b) + ε.
        m
   Notice, that
                                          m−1
                                f ′′ (ξ) X       1
                                   ′
                                        =          (n)
                                f (ξ)     i=1 ξ − xi
                                          X      1         1    X     1
                                        =          (n)
                                                       +      +
                                          i<j ξ − xi
                                                         ξ − b i>j ξ − x(n)
                                              | {z }               | {z i }
                                                      1
                                                   < ξ−a                          <0

                                                         1   1
                                       < (m − 1)           +    = 0.
                                                        ξ−a ξ−b
                                                                                  f ′′
The last equality holds by definition of ξ. Since f ′ is positive and                  is decreasing in (a, b), we have
                                                                                  f′
that f ′′ is negative on (ξ, b). Therefore,
                                          Z b           Z b
                                               ′
                          f (b) − f (ξ) =     f (t)dt ≤     f ′ (ξ)dt = (b − ξ)f ′ (ξ)
                                              ξ                      ξ

Hence,

                                 f (ξ) + f ′ (ξ) ≥ f (b) − (b − ξ)f ′ (ξ) + f ′ (ξ)
                                                 = f (b) + (1 − (ξ − b))f ′ (ξ)
                                                 = f (b) + (1 − m1 (b − a))f ′ (ξ)
                                                 ≥ f (b) + (1 − m1 )f ′ (ξ).

Together with
                                                  m−1
                         ′         ′
                                                  Y           (n)                       dm−1
                        f (ξ) = |f (ξ)| = m             |ξ − xi | ≥ m|ξ − b|m−1 ≥
                                                  i=1
                                                        |    {z      }                  mm−2
                                                            ≥|ξ−b|

we get
                                             f (ξ) + f ′ (ξ) ≥ f (b) + ε.
Together with (2) this shows (3). This finishes the proof of Lemma 2.
                                                  f + f′




                                              f


                                                                             f′
                                         a                     ξ         b


                                                              4
